\section{The partial-observation construction}
\label{sec:annex-p17-partial-observation-decisions}

\begin{annexitem}{annex:p17.astrom.filter}{thm:p17-astrom-filter}{Joint laws, action conditioning and every posterior branch}
\begin{proof}
Start with \(P(X_1=i,Y_1=j)=\mu_iq_{ij}\).
At each subsequent step multiply the current history law by
the policy kernel for the action, the specified hidden-state
transition probabilities, and the observation probabilities.
These are measurable probability kernels: the finite matrices
are stochastic and continuous in the control, and every action
kernel is Borel. Repeated kernel products, as proved in
\cref{thm:p15-kernel-joint-measure}, give a unique finite joint
law, with the declared conditional laws at each step.
For an indefinitely specified sequence the countable
standard Borel product construction in
\cref{thm:p15-coded-space-extension} gives the same finite laws.
The finite-horizon argument itself needs only the finite products.

The initial formula follows by finite Bayes division on
\(\{\lambda_j^0>0\}\). At zero probability choose the stated
default; it changes no conditional identity.
Suppose \(B_t\) is a conditional law of \(X_t\) given
\(\mathcal I_t\). Since the current control is chosen by
a kernel on \(\mathcal I_t\) without additional hidden-state
information, its joint conditional law with \(X_t\) is
\[
        P(X_t=s,U_t\in C\mid\mathcal I_t)
                        =B_t(s)\pi_t(C\mid H_t).
\]
Thus the conditional law of \(X_t\) given
\(\mathcal I_t\) and the chosen \(U_t\) is still \(B_t\).
This is an identity of kernels, valid without dividing
by the probability of a particular continuous action.

Summing over the current hidden state gives
\[
 P(X_{t+1}=i,Y_{t+1}=j\mid H_t,U_t)
          =q_{ij}\sum_s B_t(s)p_{si,t}(U_t)
          =z_i^j(B_t,U_t,t).
\]
Summing in \(i\) gives \(\lambda_j\).
For \(\lambda_j>0\), ordinary finite Bayes division gives
\(F_j(i)=z_i^j/\lambda_j\).
The probability of observing a branch with
\(\lambda_j(B_t,U_t,t)=0\) is
\[
 \sum_j E[\mathbf1_{\{\lambda_j(B_t,U_t,t)=0\}}
                                      \lambda_j(B_t,U_t,t)]=0.
\]
Consequently any fixed default on these branches gives
a version of the conditional law after the next observation.
This proves the induction, including all degenerate priors.

Each \(a_i,z_i^j,\lambda_j\) is continuous in \((b,u)\).
The positive-denominator set is open, the quotient is
continuous there, and the default gives a Borel map
on the whole domain. Also \(F_j\in\Delta_n\) and
\(\sum_j\lambda_j=\sum_i a_i=1\).
For a Borel \(C\subseteq\Delta_n\), the displayed finite
sum defining \({\cal R}_t(C\mid b,u)\) is Borel and
is a probability measure in \(C\). Its support contains
at most one point per observation, hence at most \(m\).
Repeated posteriors can reduce that number.
Zero weights contribute zero for every default, proving
the asserted invariance.
\end{proof}
\end{annexitem}

\begin{annexitem}{annex:p17.astrom.information.state}{thm:p17-astrom-information-state}{The sufficient statistic and both policy directions}
\begin{proof}
The preceding conditional law implies
\[
 E[\varphi(B_{t+1})\mid H_t,U_t]
       =\sum_j\lambda_j(B_t,U_t,t)
                                \varphi(F_j(B_t,U_t,t))
\]
for every bounded Borel \(\varphi\).
The conditional immediate cost is
\(\ell_t(B_t,U_t)\), because the conditional law of
\(X_t\) after choosing that control is \(B_t\).
These are the precise transition and cost identities
needed for a controlled information state.
If \(U_t=f_t(B_t)\), taking conditional expectations
given the whole belief history reduces the right side
to a function of \(B_t\) alone. This proves the
time-dependent Markov assertion.
For arbitrary history controls this final reduction is
not asserted: the action can depend on more than \(B_t\).

For implementation, start from \(B_1=F^0_{Y_1}\).
Given any Borel policy on belief histories, reconstruct
the successive beliefs from the observations and the
known actions by the Borel filter. Apply the belief-history
action kernel to that reconstructed history. This is
an admissible original action kernel. Induction in the
joint kernel law shows that its belief/action finite
distributions equal those of the belief model: they
have the same initial belief law, action choices, and
conditional transition \({\cal R}_t\).
The equality of expected stage costs follows from
the displayed conditional cost identity, and summation
gives equality of the total costs.

Conversely fix an original policy and its initial law.
For each \(t\), take a regular conditional law
\(\kappa_t(du\mid b)\) of \(U_t\) given \(B_t=b\).
The target \(U\), a compact Euclidean subset, is standard
Borel. Thus \cref{thm:p15-value-conditional-law} supplies
such a Borel kernel, with arbitrary measurable default
on its null exceptional set. If \(\nu_t\) is the
original law of \(B_t\), its joint law with \(U_t\)
is \(\nu_t(db)\kappa_t(du\mid b)\).

In a new belief model use \(\kappa_t\) as the Markov
action kernel. Inductively its belief marginal is
\(\nu_t\): this is true initially; equality of the
belief marginal gives equality of the current joint
\((B_t,U_t)\) law, and integrating the common
\({\cal R}_t\) gives equality of the next belief marginal.
Hence all stagewise belief/action laws and the expected
costs \(\int\ell_t\,d(\nu_t\kappa_t)\) agree.
This proves the projected-cost statement. It does not
assert equality of every original history distribution,
or one projection valid for all initial laws.
Finally, in the deterministic source class, earlier
controls are functions of the successive observation
prefixes by induction on time; substituting those functions
eliminates explicit earlier-control arguments. The filtered
deterministic policy constructed below therefore belongs
to that original class.
\end{proof}
\end{annexitem}

\begin{annexitem}{annex:p17.astrom.continuous.selection}{thm:p17-astrom-continuous-selection}{The zero-mass limit and an explicit compact minimizer}
\begin{proof}
The simplex is compact. A continuous \(v\) on it is
bounded, say by \(M\). For \(\lambda_j>0\), the function
\(\lambda_j v(z^j/\lambda_j)\) is continuous.
At a point where \(\lambda_j=0\), its absolute value
on all nearby positive branches is at most \(M\lambda_j\),
which tends to zero. Its assigned zero value therefore
extends it continuously. A finite sum is continuous.
Notice that this does not claim continuity of the
normalized posterior itself at a zero denominator.

For completeness, the compactness facts needed for
minimization can be obtained directly. Every sequence
in a compact metric space has a convergent subsequence.
If it had no cluster point, each point would have a
neighborhood containing only finitely many sequence terms,
and a finite subcover would give a contradiction.
At a cluster point choose successively later terms
within \(1/k\) to obtain the subsequence.
A continuous real function on a compact space is
bounded, since otherwise points with absolute values
tending to infinity have a convergent subsequence.
A sequence approaching its infimum then has a convergent
subsequence, and continuity gives attainment.

Let \(h(b)=\min_{u\in U}H(b,u)\).
For \(b_k\to b\), choose minimizing \(u_k\).
On any convergent subsequence \(u_{k_l}\to u\),
continuity gives
\(\lim_l h(b_{k_l})=H(b,u)\geq h(b)\).
Conversely a minimizer \(u^*\) at \(b\) gives
\(h(b_k)\leq H(b_k,u^*)\to h(b)\).
If \(h(b_k)\) did not tend to \(h(b)\), choose a
subsequence violating convergence; compactness of its
minimizers and these two inequalities give a contradiction.
Thus \(h\) is continuous. The same proof applies when
\(U\) is replaced by any nonempty compact subset \(K\).

Here is an explicit Borel minimizing selector.
Put \(U\) inside a closed bounded cube \(C_0\).
Bisect every coordinate repeatedly; at each level the
finitely many closed child cubes cover their parent,
with overlaps on boundaries allowed. For a fixed
cube \(C\) with \(U\cap C\ne\varnothing\), the function
\[
              h_C(b)=\min_{u\in U\cap C}H(b,u)
\]
is continuous by the preceding argument.
The condition \(h_C(b)=h(b)\) is closed and says
exactly that this cube contains a minimizer.
Starting at \(C_0\), choose the first child satisfying
that condition, and repeat inside the chosen child.
At least one child always qualifies because a current
minimizer belongs to a child. By induction the selected
cube index is Borel: on each set where the previous
index is fixed, the choice is the first among finitely
many closed comparison conditions.

The selected cubes are nested, their diameters tend to
zero, and every one contains a point of the compact
minimizer set. Taking the limit of these points gives
a unique point \(f(b)\in U\) with \(H(b,f(b))=h(b)\).
The centers of the selected cubes are finite-range
Borel maps converging to \(f\); each Euclidean coordinate
of their limit is Borel. Hence \(f\) is Borel.
This argument selects on the full simplex. The cubes
are a proof construction in the control space, not
a discretization of the belief state.
\end{proof}
\end{annexitem}

\begin{annexitem}{annex:p17.astrom.finite.optimum}{thm:p17-astrom-finite-optimum}{Backward existence and optimality against all histories}
\begin{proof}
At the last date \(\mathcal H_N(b,u)=\ell_N(b,u)\) is continuous.
The preceding compact-minimum lemma gives continuous
\(V_N\) and a Borel minimizing \(f_N\).
If \(V_{t+1}\) is continuous, the weighted-posterior
lemma makes its continuation term continuous in \((b,u)\).
The immediate cost is continuous because it is a
finite sum of products of coordinates and continuous
cost functions. Hence \(\mathcal H_t\) is continuous, and the
same minimum lemma gives \(V_t,f_t\).
This proves the construction by backward induction.
Its values are bounded by
\(\sum_{k=t}^N\max_{i,u}|g_k(i,u)|\).
The terminal condition and backward equations also
prove uniqueness of their sequence of solutions.

For the control argument, let an arbitrary policy be
used from a history at stage \(t\), and suppose its
future conditional cost from stage \(t+1\) is bounded
below by \(V_{t+1}(B_{t+1})\).
The filter and conditional cost identities then bound
its current conditional cost, after choosing action \(u\),
below by \(\mathcal H_t(B_t,u)\).
Every \(\mathcal H_t(B_t,u)\geq V_t(B_t)\); averaging against
any randomized history action preserves this inequality.
The base at stage \(N\) is precisely the minimum of
the last immediate cost. This proves by backward
induction the lower bound for all policies.

For \(U_t=f_t(B_t)\), equality holds in the minimum
at every \(b\), and the induction has equality at each
subsequent step. Thus its conditional cost is \(V_t(b)\)
from every starting belief and date.
Initial averaging gives
\(\sum_j\lambda_j^0V_1(F_j^0)\).
Implementation in the original model is supplied by
the preceding information-state proof. Conversely,
every original policy is bounded below by the same
value; or its cost can be transferred to the belief
model by the Markov projection proved there.
Therefore the two models have the same minimum and
both possess the constructed deterministic optimum.
The necessity and sufficiency asserted by the source's
Theorems 1 and 2 follow from this equality and the
verification identity next proved, with the displayed
terminal condition retained.
\end{proof}
\end{annexitem}

\begin{annexitem}{annex:p17.astrom.verification}{thm:p17-astrom-verification}{The residual identity and conditional future costs}
\begin{proof}
The filter theorem gives
\[
 E[V_{t+1}(B_{t+1})\mid H_t,U_t]
                  =\mathcal H_t(B_t,U_t)-\ell_t(B_t,U_t).
\]
Take expectations and use
\(E\ell_t(B_t,U_t)=E g_t(X_t,U_t)\).
Rearrangement gives
\[
 E\delta_t(B_t,U_t)
   =E g_t(X_t,U_t)+E V_{t+1}(B_{t+1})-E V_t(B_t).
\]
For \(t=N\), use the declared zero terminal value.
Summing makes the intermediate value terms telescope,
proving the identity.
Every residual is nonnegative and bounded, so their
sum of expectations is zero exactly when each residual
is zero almost surely. This proves necessity and
sufficiency for the specified initial law.
If all residuals are bounded by \(\eta_t\), their
expectations give the stated total-loss bound.

To state the history-dependent value correctly, put
\[
 W_t^\pi(H_t)=
              E^\pi[\sum_{k=t}^N g_k(X_k,U_k)\mid H_t].
\]
It obeys the conditional recursion
\[
 W_t^\pi(H_t)=
  \int_U\left[\ell_t(B_t,u)+
      \sum_j\lambda_j(B_t,u,t)
        W_{t+1}^\pi(H_t,u,j)\right]\pi_t(du\mid H_t),
\]
with zero future term at \(N\).
Here \(H_t,u,j\) is the extended observation/action history.
The expression is justified by finite kernel integration
and conditional expectation, including chosen versions
at null histories. For an optimal belief policy,
\(W_t^\pi=V_t(B_t)\) almost surely.
The terms involving future \(B_k\) are random before
those future observations occur, so their unconditioned
sum cannot in general be identified with \(W_t^\pi\)
at the current history.
\end{proof}
\end{annexitem}

\begin{annexitem}{annex:p17.astrom.information.bounds}{thm:p17-astrom-information-bounds}{Every stage of both information comparisons}
\begin{proof}
The perfect-information recursion follows by conditioning
on the observed current state \(i\), minimizing over
the compact control set, and using the next-state
probabilities. The same finite-horizon induction as
above proves its value and attainment.
For a current distribution \(b\), revealing the current
state first gives value \(L_t(b)=\sum_i b_iS_t(i)\).

We prove \(L_t\leq V_t\) by induction.
At \(N\), for every \(u\),
\(\sum_i b_i g_N(i,u)\geq\sum_i b_i\min_vg_N(i,v)\);
minimize the left side.
At an earlier date, the induction hypothesis and
the identity
\(\sum_j\lambda_jF_j(i)=\sum_s b_sp_{si,t}(u)\)
give
\[
 \begin{split}
 \mathcal H_t(b,u)
 &\geq \sum_s b_s
       [g_t(s,u)+\sum_i p_{si,t}(u)S_{t+1}(i)]\\
 &\geq \sum_s b_s S_t(s)=L_t(b).
 \end{split}
\]
Zero-weight branches satisfy the same identity because
their unnormalized vectors are zero. Minimization
proves the next induction step.

For no further observations, a fixed control schedule
\((u_t,\ldots,u_N)\) gives deterministic predicted laws
\(b_t=b\), \(b_{k+1}=b_kP_k(u_k)\), and expected cost
\(\sum_{k=t}^N\ell_k(b_k,u_k)\).
This is continuous on the compact finite product
\(U^{N-t+1}\), so its minimum is attained.
Splitting off the first control proves exactly the
displayed \(O_t\) recursion. In the original problem
this schedule is an admissible policy which ignores
new observations. Thus \(V_t(b)\leq O_t(b)\).
This restriction-of-policies argument supplies the
source's upper comparison without replacing the
time-ordered products by a single matrix power.

If \(q_{ij}=r_j\) for every \(i\), then
\(z_i^j=r_ja_i\), \(\lambda_j=r_j\), and every
positive-probability posterior is \(a=bP_t(u)\).
The Bellman continuation is therefore
\(\sum_jr_j V_{t+1}(a)=V_{t+1}(a)\), including
zero \(r_j\)'s. The recursion equals that for \(O_t\).
A deterministic optimum is consequently a control
schedule once the initial belief is fixed.
This does not assert that every feedback policy has
deterministic beliefs: its controls may use the
uninformative measurements as randomization.
If \(Q\) is the identity, observation \(j\) of positive
probability has posterior \(\delta_j\).
The belief recursion restricted to these point masses
is exactly \(S_t\), proving the perfect-observation case.

For the initial comparison, average \(L_1\leq V_1\leq O_1\)
under the initial posterior law.
Bayes' numerator identity gives
\(\sum_j\lambda_j^0F_j^0(i)=\mu_i\), so the averaged
lower bound is \(\sum_i\mu_iS_1(i)\).
For any fixed open-loop schedule, its expected cost
is linear in its starting distribution.
Its \(\mu\)-cost thus equals the average of its
\(F_j^0\)-costs. On each posterior its cost is at
least \(O_1(F_j^0)\). Minimize over the original
schedule to obtain
\(\sum_j\lambda_j^0O_1(F_j^0)\leq O_1(\mu)\).
This proves the full four-term comparison and
identifies exactly which first observation is available.
Subtracting the two measured-information values from
the same no-observation baseline proves all the
assertions about nonnegative information values.
\end{proof}
\end{annexitem}

\begin{annexitem}{annex:p17.astrom.delays}{thm:p17-astrom-delays}{A finite augmented state for delayed observations}
\begin{proof}
At date \(t\), use hidden state
\(Z_t=(X_1,\ldots,X_t)\).
There are finitely many such paths. The next transition
appends \(i\) with probability \(p_{X_ti,t}(u)\);
all other next-path probabilities are zero.
These transition entries are continuous in \(u\).
The current cost is \(g_t(X_t,u)\), continuous in \(u\).
One may use a single finite set consisting of all
paths of lengths at most \(N\), with length tags;
transitions at date \(t\) use only the corresponding
length-\(t\) part. This avoids any change of the
finite-state hypothesis.

The new observation delivered at date \(t+1\) is the
finite block of readings indexed
\(d_t+1,\ldots,d_{t+1}\).
For a hidden path \(z=(x_1,\ldots,x_{t+1})\), its
probability of being the block \((j_r)\) is
\(\prod_{r=d_t+1}^{d_{t+1}}q_{x_rj_r}\).
An empty product equals one and has a singleton
observation alphabet. Independent observation generation
with the stated conditional laws justifies this product:
these readings were not previously delivered, and
could not influence earlier controls before delivery.
The initial delivered block has the analogous likelihood.
The emission matrix now depends on time and the
finite augmented state. Every row sums to one by
summing the product of its finite probabilities.

Apply the already proved filter proof with these
finite path states and finite block observations.
All its sums and conditional identities remain
literal finite sums; using a time-indexed observation
matrix only replaces each \(q_{ij}\) by the indicated
date's fixed entry. Its posterior lives on the full
simplex of the augmented finite state.
Each transition entry is continuous in \(u\), and
the same weighted-zero-mass continuity proof applies.
The compact-minimizer and backward-optimality proofs
therefore give an attained deterministic policy for
this delayed information problem.
The controller can implement it by updating the
path posterior from newly delivered readings and
its known actions.

If a delivery rule supplies at least the readings
of another at every date, a controller may ignore
the extra readings and reproduce every policy of
the smaller-information problem. Its optimum is
therefore no larger. With \(d_t=t\) the current
reading is available, and the original current-state
filter is sufficient. With \(d_t=0\) no readings
arrive and the attained optimum is the schedule
already constructed.
The nondecreasing delivery frontier records retention
of previously received readings. A decreasing imposed
memory cutoff describes forced forgetting and is
not asserted to have this same information structure.
\end{proof}
\end{annexitem}

\begin{annexitem}{annex:p17.astrom.example.one}{thm:p17-astrom-example-one}{All four and nine branches of the switch calculation}
\begin{proof}
The predicted probability of state one is
\(a=(1-u)x+u(1-x)\). As \(u\) ranges over \([0,1]\),
this ranges over the closed interval with endpoints
\(x,1-x\); when \(x=1/2\) it is constant.
Put
\[
 A=qa,\quad B=(1-q)(1-a),\quad
 C=(1-q)a,\quad D=q(1-a).
\]
These are, in order, \(z_1^1,z_2^1,z_1^2,z_2^2\).
Their column sums are the observation probabilities.
This proves the posterior formulas in the definition;
the zero-column cases have zero probability and zero
contribution to each calculation that follows.

The terminal cost gives \(V_4(x)=1-x\).
At stage three, the expected terminal cost is
\(B+D=1-a\), so minimizing over the stated interval
gives \(V_3(x)=\min(x,1-x)\).
If \(x<1/2\), maximizing \(a\) takes \(u=1\);
if \(x>1/2\), it takes \(u=0\); at equality every
control works.

For stage two, weighting the next stage's normalized
minimum by its observation probability cancels the
normalization. Thus the expression to minimize is
\[
                   \Phi_2(a)=\min(A,B)+\min(C,D).
\]
The four branches, including their exact consistency
conditions, are
\[
 \begin{array}{c|c|c}
 \text{chosen terms}&\text{sum}&\text{conditions}\\ \hline
 A,C&a&A\leq B,\ C\leq D\\
 A,D&q&A\leq B,\ D\leq C\\
 B,C&1-q&B\leq A,\ C\leq D\\
 B,D&1-a&B\leq A,\ D\leq C .
 \end{array}
\]
The conditions are respectively
\(a\leq e\), \(q\leq a\leq1-q\),
\(1-q\leq a\leq q\), and \(a\geq1-e\).
The middle branches therefore require, respectively,
\(q\leq1/2\) and \(q\geq1/2\).
Equality can satisfy several rows; the row sums then
agree. Equivalently, the sum of the two minima is
the minimum of all four sums, without needing to
exclude nonselected branches separately:
\[
                      \Phi_2(a)=\min(a,1-a,e).
\]
This is a concave function, because for affine \(L_i\)
one has
\(\min_iL_i(\theta a+(1-\theta)b)
 \geq\theta\min_iL_i(a)+(1-\theta)\min_iL_i(b)\).
It is symmetric under \(a\mapsto1-a\).
A concave function on an interval is at least the
smaller of its endpoint values; here the two endpoint
values agree. Thus its minimum over the predicted
interval is \(\Phi_2(x)\), attained at \(u=0\) and \(u=1\).
If \(e\leq x\leq1-e\), that whole interval lies in
the constant part \(\Phi_2=e\), so every \(u\) is optimal.

For stage one the expression is
\[
 \Phi_1(a)=\min(A,B,e(A+B))+\min(C,D,e(C+D)).
\]
Here are all nine branches. A branch is selected
exactly when its left summand is no larger than the
other two left candidates and its right summand is
no larger than the other two right candidates.
These non-strict comparisons include every tie boundary.
\[
 \begin{array}{c|ccc}
  &C&D&e(C+D)\\ \hline
 A&A+C&A+D&A+e(C+D)\\
 B&B+C&B+D&B+e(C+D)\\
 e(A+B)&e(A+B)+C&e(A+B)+D&e .
 \end{array}
\]
The minimum of the nine displayed affine sums equals
the sum of the two minima, so this finite table also
gives an exact expression on every equality boundary.
Interchanging the two observation names replaces \(q\)
by \(1-q\) and leaves this sum unchanged.
It therefore suffices to simplify it for \(q=1-e\),
with \(0\leq e\leq1/2\).
Let
\[
 L(a)=\zeta+(\alpha-\zeta)a,\qquad
 K=1-2e+2e^2,\qquad c=e(1-e).
\]
In row order, the nine sums become
\[
 \begin{array}{ccc}
 a&1-e&Ka+c\\
 e&1-a&L(1-a)\\
 L(a)&K(1-a)+c&e .
 \end{array}
\]
For example
\(e(A+B)+C=e^2+(2e-2e^2)a=L(a)\);
direct substitution gives the other eight entries.
The extra third entry satisfies
\[
           Ka+c-L(a)=e(1-2e)+(1-2e)^2a\geq0,
\]
and its reflected entry is similarly no smaller
than \(L(1-a)\).
Moreover
\(\min(L(a),L(1-a))\leq(L(a)+L(1-a))/2=e\),
and \(1-e\geq e\).
Thus the lower envelope reduces exactly to
\[
                \Phi_1(a)=\min(a,1-a,L(a),L(1-a)).
\]
It is again concave and symmetric, so its minimum
over the predicted interval is \(\Phi_1(x)\), with
both endpoint controls optimal. This proves \(V_1\)
and accounts for every original branch.
At \(e=0\) the zero lines give \(V_2=V_1=0\);
at \(e=1/2\) the formulas reduce to
\(\min(x,1-x)\). Hence perfect, reversed-perfect,
uninformative and degenerate-belief endpoints are included.

For \(q=4/5\), one has \(e=1/5\),
\(\alpha=9/25,\zeta=1/25\).
At \(x=1/10\),
\[
 L(x)=\frac9{25}\frac1{10}
       +\frac1{25}\frac9{10}
       =\frac9{125}<\frac1{10},
\]
giving all four stated values.
One concrete optimal implementation chooses \(u=0\)
at stages one and two, then uses the stage-three
posterior to flip the more likely state two to state one.
For two additional readings with accuracy \(4/5\),
at prior \(1/10\) the final decision regards state
one as more likely only when both readings say one.
Its error probability is
\[
       \frac1{10}(1-(4/5)^2)
              +\frac9{10}(1/5)^2
                       =\frac9{250}+\frac9{250}=\frac9{125}.
\]
For example, after the first reading says one its
probability is \(13/50\) and the posterior is \(4/13\);
if the next says one it becomes \(16/25\).
The other two-reading posteriors are \(1/10,1/10,1/145\),
so the policy just described is exactly the threshold
rule, not an appeal to an unobserved state.

With perfect information at a nonterminal stage,
choose \(u=0\) in state one and \(u=1\) in state two,
and keep state one thereafter. The terminal failure
probability is zero, giving \(L_1=L_2=L_3=0\);
at the terminal stage no control changes the accrued
state cost, giving \(L_4=1-x\).
For the schedule comparison, \(O_3=\min(x,1-x)\).
Its preceding recursion minimizes
\(\min(a,1-a)\) over the same interval, which by the
concavity and symmetry argument remains \(\min(x,1-x)\).
Induction gives \(O_2,O_1\), and \(O_4=1-x\).
\end{proof}
\end{annexitem}

\begin{annexitem}{annex:p17.astrom.example.two}{thm:p17-astrom-example-two}{All action values in the profit calculation}
\begin{proof}
The immediate expected reward of action \(u\) is
\(25x-c_u\), with
\((c_1,c_2,c_3,c_4)=(5,8,15,18)\).
The transition table gives the four predicted
probabilities \(a_u\) in the definition.
For completeness the four unnormalized posterior
columns, in the order \(z_1^1,z_2^1,z_1^2,z_2^2\), are
\[
 \begin{array}{c|rrrr}
 u&z_1^1&z_2^1&z_1^2&z_2^2\\ \hline
 1&.32+.08x&.18-.03x&.08+.02x&.42-.07x\\
 2&.56-.16x&.09+.06x&.14-.04x&.21+.14x\\
 3&.32+.32x&.18-.12x&.08+.08x&.42-.28x\\
 4&.56+.08x&.09-.03x&.14+.02x&.21-.07x .
 \end{array}
\]
They follow respectively from
\(.8a_u,.3(1-a_u),.2a_u,.7(1-a_u)\).
Every \(a_u\) lies in \([.4,.8]\);
thus \(\lambda_1=.3+.5a_u\) and
\(\lambda_2=.7-.5a_u\) are strictly positive.
The rational posterior formulas are therefore valid
throughout this example.

At stage four the largest immediate reward is action
one's \(25x-5\), strictly above the other three.
At stage three its next-stage expectation is
\(25a_u-5\); adding the immediate reward gives
all four candidate functions
\[
 \begin{array}{c|r}
 u&\mathcal H_3(x,u)\\ \hline
 1&27.5x\\
 2&20x+4.5\\
 3&35x-10\\
 4&27.5x-5.5 .
 \end{array}
\]
Action three is below action one, since their
difference is \(10-7.5x\geq2.5\).
Action four is below action one by \(5.5\).
The remaining two intersect at \(x=3/5\), with
action two better below and action one better above.
This proves \(V_3\) and its complete action rule.

For stage two, the high branch of \(V_3\) is
\(27.5b\) and the low branch is \(20b+4.5\).
After observation one the posterior exceeds \(3/5\)
exactly when \(a_u>9/25=.36\).
After observation two it exceeds \(3/5\) exactly
when \(a_u>21/25=.84\).
The entire range \([.4,.8]\) lies strictly between
these thresholds. Thus observation one always selects
the high branch and observation two the low branch.
Their weighted sum is exactly
\[
 27.5(.8a_u)+20(.2a_u)+4.5(.7-.5a_u)
                                      =23.75a_u+3.15.
\]
Adding \(25x-c_u\) gives every candidate:
\[
 \begin{array}{c|r}
 u&\mathcal H_2(x,u)\\ \hline
 1&27.375x+7.650\\
 2&20.250x+11.775\\
 3&34.500x-2.350\\
 4&27.375x+1.775 .
 \end{array}
\]
The first minus the third is \(10-7.125x\geq2.875\);
the first minus the fourth is \(5.875\).
The first two intersect at
\((11.775-7.650)/(27.375-20.250)=11/19\).
This proves \(V_2\) and its policy, with both actions
permitted at the exact rational threshold.

At stage one the high branch of \(V_2\) is
\(27.375b+7.650\), and the low branch is
\(20.250b+11.775\), separated at \(11/19\).
The observation-one posterior is above that threshold
exactly when \(a_u>33/97\); the observation-two
posterior is above it exactly when \(a_u>77/93\).
Again \([.4,.8]\) lies strictly between the two.
Hence the weighted next value is
\[
 \begin{split}
 &27.375(.8a_u)+7.650(.3+.5a_u)\\
 &\qquad+20.250(.2a_u)+11.775(.7-.5a_u)
                         =23.8875a_u+10.5375.
 \end{split}
\]
The four stage-one candidates are consequently
\[
 \begin{array}{c|r}
 u&\mathcal H_1(x,u)\\ \hline
 1&27.38875x+15.09250\\
 2&20.22250x+19.25875\\
 3&34.55500x+5.09250\\
 4&27.38875x+9.25875 .
 \end{array}
\]
The first minus the third is
\(10-7.16625x\geq2.83375\), and the first minus
the fourth is \(5.83375\). The first two intersect at
\[
       \frac{19.25875-15.09250}{27.38875-20.22250}
                       =\frac{3333}{5733}=\frac{1111}{1911}.
\]
All stated values, strict dominations and endpoint
policy rules now follow.
Only arithmetic on the original rational inputs has
been used. The decimal expressions and rational thresholds
are exact.
\end{proof}
\end{annexitem}

\begin{annexitem}{annex:p17.astrom.example.two.bounds}{thm:p17-astrom-example-two-bounds}{All full-information and schedule action tables}
\begin{proof}
For full information, the terminal state-value vector
is \((20,-5)\), from action one in both states.
The next three backward action tables have rows
for states one and two and columns for controls
one through four:
\[
 \begin{array}{c|rrrr}
  &1&2&3&4\\ \hline
 t=3,\ X=1&27.5&24.5&25&22\\
 t=3,\ X=2&0&4.5&-10&-5.5\\ \hline
 t=2,\ X=1&36&33&32.9&29.9\\
 t=2,\ X=2&8.7&12.6&-1.3&2.6\\ \hline
 t=1,\ X=1&44.30&41.30&41.32&38.32\\
 t=1,\ X=2&16.96&20.98&6.96&10.98 .
 \end{array}
\]
For example the stage-two entry at state two, control
two is \(-8+.7(27.5)+.3(4.5)=12.6\).
Every entry is obtained by adding the current reward
to its transition row times the next state's optimum.
The statewise maxima are, successively,
\((27.5,4.5)\), \((36,12.6)\), and \((44.3,20.98)\).
Taking their expectations under \((x,1-x)\) proves
all the \(L_t\) lines in the statement.

Without future observations, the predicted belief
is \(a_u\) deterministically. At stage three, the
candidate table is the same four functions already
computed, so \(O_3=\max(27.5x,20x+4.5)\).
At stage two the four candidates, retaining both
possible next branches where needed, are
\[
 \begin{array}{c|l}
 u&\mathcal H_2^{\rm open}(x,u)\\ \hline
 1&27x+7.5\\
 2&\max(19.5x+11.25,\ 21x+10.5)\\
 3&\max(36x-4,\ 33x-2.5)\\
 4&27.75x+1.25 .
 \end{array}
\]
These follow by substituting \(a_u\) into \(O_3\)
and adding \(25x-c_u\).
The two lines of action three are below action one's
line: the differences are \(11.5-9x\geq2.5\)
and \(10-6x\geq4\). Action four is also below
action one, by \(6.25-.75x\geq5.5\).
The extra action-two line \(21x+10.5\) is below
\(19.5x+11.25\) for \(x\leq1/2\), and below
\(27x+7.5\) for \(x\geq1/2\).
Thus
\[
                 O_2(x)=\max(27x+7.5,\ 19.5x+11.25),
\]
and its exact switching threshold is \(1/2\).

At stage one, substituting the same four \(a_u\)
into this last two-branch expression gives
\[
 \begin{array}{c|l}
 u&\mathcal H_1^{\rm open}(x,u)\\ \hline
 1&26.95x+14.05\\
 2&19.6x+18.4\\
 3&\max(35.8x+3.3,\ 32.8x+4.05)\\
 4&27.7x+8.4 .
 \end{array}
\]
For action one, \(a_1=.4+.1x\leq.5\), so the
low branch of \(O_2\) gives
\(25x-5+19.5(.4+.1x)+11.25=26.95x+14.05\).
For action two, \(a_2=.7-.2x\geq.5\), so the
high branch gives \(19.6x+18.4\).
Action four also uses the high branch. Action three
can use either branch, giving the two displayed lines.

Both action-three lines are below action one's line,
by \(10.75-8.85x\geq1.9\) and
\(10-5.85x\geq4.15\), respectively.
Action four is below action one by
\(5.65-.75x\geq4.9\).
The two remaining lines intersect at
\[
                  \frac{18.4-14.05}{26.95-19.6}
                       =\frac{4.35}{7.35}=\frac{29}{49}.
\]
This proves the final open-loop value and threshold.
At the terminal date \(O_4=25x-5\), completing
the four-stage schedule calculation.

All information bounds reverse under reward maximization,
because replacing rewards by their negatives transforms
the problem into the proved cost-minimization comparison.
Thus \(O_t\leq V_t\leq L_t\).
At \(x=1/2\), direct evaluation gives
\[
 O_1=19.6/2+18.4=28.2,\quad
 V_1=20.2225/2+19.25875=29.37,\quad
 L_1=23.32/2+20.98=32.64 .
 \]
Subtracting the common schedule baseline yields
\(1.17\) and \(4.44\). These are values conditional
on the declared current belief. They do not silently
change which initial observation was already received.
\end{proof}
\end{annexitem}

\begin{annexitem}{annex:p17.astrom.information.failures}{thm:p17-astrom-information-failures}{Five counterexamples and a fixed-parameter augmentation}
\begin{proof}
For a zero-likelihood example, let the hidden state
be any fixed state and let observation one have
probability one at every state. Observation two
then has numerator vector zero and denominator zero.
Every default probability vector on that event
defines the same conditional-law version on all
positive-probability events. Division by its
likelihood would be \(0/0\), which defines no vector.
The filter's weighted contribution on that event
is zero for every default.

For temporal dependence, use a fixed hidden fair bit
\(X\in\{0,1\}\). Let \(E\) be an independent bit with
\(P(E=0)=3/4\), and set \(Y_1=Y_2=X\mathbin{\mathrm{xor}}E\).
At each date the one-time observation matrix has
accuracy \(3/4\). But two observations both equal to
one carry exactly the same information as the first,
so the true posterior \(P(X=1\mid Y_1=Y_2=1)\) is
\(3/4\). The independent-observation update would
instead give
\((3/4)^2/[(3/4)^2+(1/4)^2]=9/10\).
Thus the one-time observation matrix alone is
insufficient; the model's conditional observation
law given the past is a substantive hypothesis.

Even mutual independence of named sensor noises does
not specify their relationship with the hidden state.
Let \(E_1,E_2\) be independent fair bits, let the
constant hidden state be \(X=E_1\mathbin{\mathrm{xor}}E_2\),
and put \(Y_t=X\mathbin{\mathrm{xor}}E_t\) for \(t=1,2\).
Then \(Y_1=E_2\), \(Y_2=E_1\), and every one-time
observation row is fair. One reading leaves \(X\)
fair, whereas the two readings determine
\(X=Y_1\mathbin{\mathrm{xor}}Y_2\).
The row-wise independent filter would leave probability
\(1/2\) after both readings. The explicit conditional
observation law rules out this coupling as well.

For observed costs, keep the hidden bit fixed and
let the stated observations be constant.
At stage one incur the state-dependent cost \(X\);
at stage two choose a guess \(u\in\{0,1\}\) and
incur \(\mathbf1_{\{u\ne X\}}\).
If the first realized cost is observed, it reveals
\(X\) and permits second-stage cost zero.
The total expected cost is then \(1/2\).
If that extra signal is omitted, the belief remains
fair and the smallest expected second-stage cost is
\(1/2\), giving total cost one.
Therefore observed costs must be represented in the
actual information history when they convey data.

For uninformative feedback, let \(X_1=0\) be known,
let every \(Y_t\) be an independent fair bit, and
let the transition after choosing \(u\in\{0,1\}\)
set \(X_{t+1}=u\) deterministically.
Take \(U_1=Y_1\). The observation matrix has identical
rows, yet \(B_2=\delta_{U_1}=\delta_{Y_1}\).
The belief depends on the observed bit through
the known action and its controlled transition.
Conditional on that action the next uninformative
reading changes nothing, exactly as the filter proves.

For an unreachable belief, take a one-stage hidden
bit known initially to be zero and an uninformative
observation. Guessing zero is optimal at that actual
belief. Extend the policy to choose zero also at
belief \(\delta_1\); there it has loss one although
guessing one has loss zero. The policy is still
optimal for its specified initial law.
This proves why a necessary optimality condition
cannot be imposed at all unreachable beliefs.

Finally suppose a parameter \(\Theta\) belongs to a
finite set and is drawn initially but remains fixed.
If the hidden transition given \(\theta\) is
\(p_{si,t}^\theta(u)\), define the augmented finite
state transition by
\[
 P((X_{t+1},\Theta)=(i,\theta')\mid(X_t,\Theta)=(s,\theta),u)
                  =\mathbf1_{\{\theta'=\theta\}}
                                      p_{si,t}^\theta(u).
\]
Its rows sum to one, and continuity in \(u\) is
preserved. Use the specified observation likelihood
for \((i,\theta)\) and the specified joint initial law.
The proved finite filter and control construction
then apply to the joint posterior.

For a concrete failure of averaging only once,
let \(\Theta\) be equally likely \(1/5\) or \(4/5\).
Conditional on \(\Theta\), successive visible bits
are independent Bernoulli with parameter \(\Theta\).
After a first success, Bayes' rule gives posterior
weights \(1/5,4/5\) on the two parameters.
The next success probability is
\[
                   (1/5)(1/5)+(4/5)(4/5)=17/25.
\]
The transition averaged once at the initial prior
would predict \(1/2\) again. It therefore fails to
represent the updated prediction, whereas the
augmented-state filter gives \(17/25\).
\end{proof}
\end{annexitem}
